%% toc: P3 34 34.1,34.2,34.3,34.4
\section{Proofs: independence, simplicity and their intersections}
\label{sec:annex-p03-w03-p03-035}

Throughout, \(\mathfrak C\) is the \(\bar\kappa\)-saturated model of
\cref{def:p03-w03-p03-035-ambient}, with \(\bar\kappa>2^{\abs T}\).  ``Realized''
means realized in \(\mathfrak C\).  Two facts about \(\mathfrak C\) are used
repeatedly, both from \cref{thm:p03-w03-p03-031-saturation}: \emph{(Sat)} a
set of formulas in a small number of variables, with small parameter set,
every finite subset of which is realized, is realized; \emph{(Copy)} if
\(c\equiv_A c'\) and \(e\) is a small tuple, there is \(e'\) with
\(ce\equiv_A c'e'\).  Since \(T\) is complete, a sentence true in one model of
\(T\) is true in \(\mathfrak C\).

\begin{annexitem}{annex:p03-w03-p03-035-em-extraction}{thm:p03-w03-p03-035-em-extraction}{Indiscernible sequences realizing an Ehrenfeucht--Mostowski type}
\begin{proof}
Let \(E\) be the set of \(\Lang(A)\)-formulas \(\theta(z_1,\ldots,z_n)\),
\(n\geq1\), true of every increasing \(n\)-tuple of the sequence
\((a_i:i<\omega)\).  In variables \((x_t:t\in X)\), each of the length of the
\(a_i\), let \(\Sigma\) consist of
\begin{enumerate}
  \item \(\theta(x_{t_1},\ldots,x_{t_n})\) for \(\theta\in E\) and
        \(t_1<\cdots<t_n\) in \(X\); and
  \item \(\psi(x_{t_1},\ldots,x_{t_n})\leftrightarrow\psi(x_{s_1},\ldots,x_{s_n})\)
        for every \(\Lang(A)\)-formula \(\psi(z_1,\ldots,z_n)\) and all
        \(t_1<\cdots<t_n\) and \(s_1<\cdots<s_n\) in \(X\).
\end{enumerate}
A realization of \(\Sigma\) is an \(A\)-indiscernible sequence with the
required property, by the definition of order indiscernibility over \(A\)
(\cref{def:p03-order-indiscernibles}).  The variables are \(\abs X\)-many
tuples of fixed finite length and the parameters lie in \(A\), so by (Sat) it
suffices to realize each finite \(\Sigma_0\subseteq\Sigma\).

Let \(t_1<\cdots<t_m\) list the indices occurring in \(\Sigma_0\), and let
\(\psi_1,\ldots,\psi_q\) be the formulas occurring in clauses of kind~(2) of
\(\Sigma_0\), with \(\psi_j\) in \(n_j\leq m\) variables.  Put
\(H_0=\omega\).  For \(j=1,\ldots,q\), divide the \(n_j\)-element subsets
\(\{i_1<\cdots<i_{n_j}\}\) of \(H_{j-1}\) into two classes according to
whether \(\mathfrak C\entails\psi_j(a_{i_1},\ldots,a_{i_{n_j}})\).  By
Ramsey's theorem (\cref{thm:p03-w03-p03-028-ramsey-a}, with
\(\Gamma=H_{j-1}\), \(r=n_j\), \(\mu=2\)) there is an infinite
\(H_j\subseteq H_{j-1}\) all of whose \(n_j\)-element subsets lie in one
class.  Let \(h_1<\cdots<h_m\) be the first \(m\) elements of \(H_q\) and
assign \(a_{h_l}\) to \(x_{t_l}\).  A clause of kind~(1) holds because
increasing indices in \(X\) go to increasing indices in \(\omega\).  A clause
of kind~(2) for \(\psi_j\) holds because both sides are evaluated at
increasing tuples from \(H_q\subseteq H_j\), on which \(\psi_j\) has constant
truth value.  So \(\Sigma_0\) is realized.
\end{proof}
\end{annexitem}

\begin{annexitem}{annex:p03-w03-p03-035-ip-characterization}{thm:p03-w03-p03-035-ip-characterization}{Equivalent forms of the independence property}
\begin{proof}
For \(n<\omega\) let \(\sigma_n\) be the sentence
\[
  \exists x_1\cdots\exists x_n\Bigl(\bigwedge_{1\leq i<j\leq n}x_i\neq x_j
  \;\wedge\;\bigwedge_{J\subseteq\{1,\ldots,n\}}\exists y
  \bigl(\bigwedge_{j\in J}\varphi(x_j;y)\wedge
  \bigwedge_{j\notin J}\neg\varphi(x_j;y)\bigr)\Bigr),
\]
where \(x_i\neq x_j\) abbreviates the disjunction of the coordinatewise
inequalities.  A model satisfies \(\sigma_n\) exactly when it contains a set
of \(n\) distinct \(\abs x\)-tuples shattered by \(\varphi\).  A subset of a
shattered set is shattered.

\emph{(1)\(\Leftrightarrow\)(2).}  If (1) holds, then for every \(n\) some
model of \(T\) contains a shattered set of at least \(n\) tuples, hence one of
exactly \(n\) tuples, so it satisfies \(\sigma_n\); by completeness
\(\sigma_n\in T\) and \(\mathfrak C\entails\sigma_n\).  Conversely
\(\mathfrak C\) is a model of \(T\).

\emph{(2)\(\Rightarrow\)(3).}  Let \(\Sigma(x_i:i<\omega)\) consist of
\(x_i\neq x_j\) for \(i<j\) and, for all finite disjoint \(I,J\subseteq\omega\),
\(\exists y(\bigwedge_{i\in I}\varphi(x_i;y)\wedge\bigwedge_{j\in J}\neg\varphi(x_j;y))\).
A finite subset mentions only indices below some \(n\); enumerate a
shattered set of \(n\) tuples as \(a_0,\ldots,a_{n-1}\).  For disjoint
\(I,J\subseteq n\), a parameter cutting out \(\{a_i:i\in I\}\) satisfies
\(\varphi(a_i;y)\) for \(i\in I\) and \(\neg\varphi(a_j;y)\) for \(j\in J\).
So every finite subset is realized, and \(\Sigma\) is realized by (Sat),
since it has countably many variables and no parameters.

\emph{(3)\(\Rightarrow\)(2).}  For \(n<\omega\) and \(A_0\subseteq
\{a_0,\ldots,a_{n-1}\}\), apply (3) with \(I\) the indices of members of
\(A_0\) and \(J\) the other indices below \(n\).

\emph{The supplement.}  Under (3), for \(A_0\subseteq\{a_i:i<\omega\}\) the
set \(\{\varphi(a_i;y):a_i\in A_0\}\cup\{\neg\varphi(a_j;y):a_j\notin A_0\}\)
has countably many parameters and is finitely realized by (3), so it is
realized by (Sat).

\emph{(3)\(\Rightarrow\)(4).}  Fix \(n\).  Using (3), enumerate \(2^n\) of
the \(a_i\) as \((a_S:S\subseteq n)\).  For \(j<n\) apply (3) with
\(I=\{S:j\in S\}\) and \(J=\{S:j\notin S\}\), read as sets of indices, to get
\(b_j\) with \(\varphi(a_S;b_j)\) iff \(j\in S\).  For \(j\neq j'\),
\(\varphi(a_{\{j\}};b_j)\) holds and \(\varphi(a_{\{j\}};b_{j'})\) fails, so
the \(b_j\) are distinct; and for \(J\subseteq n\),
\(\varphi^{\mathrm{opp}}(b_j;a_J)\) holds iff \(j\in J\).  So
\(\{b_0,\ldots,b_{n-1}\}\) is shattered by \(\varphi^{\mathrm{opp}}\).  As
\(n\) was arbitrary, (2) holds for \(\varphi^{\mathrm{opp}}\), hence (1) for
\(\varphi^{\mathrm{opp}}\) by the equivalence already proved.

\emph{(4)\(\Rightarrow\)(3).}  The opposite of \(\varphi^{\mathrm{opp}}\) is
\(\varphi\).  If (4) holds, then (1), hence (2) and (3), hold for
\(\varphi^{\mathrm{opp}}\), and the previous paragraph applied to
\(\varphi^{\mathrm{opp}}\) gives (1), hence (3), for \(\varphi\).

\emph{(4)\(\Leftrightarrow\)(6).}  Let \(B\) be a set of \(n\) distinct
\(\abs y\)-tuples.  A \(\varphi\)-type \(p\) over \(B\) contains exactly one of
\(\varphi(x;b)\), \(\neg\varphi(x;b)\) for each \(b\in B\): it contains at
most one since it is realized, and at least one by maximality, since a
realization satisfies one of them.  So \(p\mapsto
B_p=\{b\in B:\varphi(x;b)\in p\}\) is injective.  A set \(B_0\subseteq B\) is
some \(B_p\) iff some \(a\) satisfies \(\varphi(a;b)\Leftrightarrow b\in B_0\)
for \(b\in B\), that is, iff \(\varphi^{\mathrm{opp}}(\mathfrak C;a)\cap
B=B_0\) for some \(a\).  Hence \(\abs{S_\varphi(B)}\leq2^n\), with equality
iff \(B\) is shattered by \(\varphi^{\mathrm{opp}}\).  Thus (6) says that
(2) holds for \(\varphi^{\mathrm{opp}}\), which is (4).

\emph{(3)\(\Rightarrow\)(5).}  Apply
\cref{thm:p03-w03-p03-035-em-extraction} to \((a_i:i<\omega)\) with
\(A=\varnothing\) and \(X=\omega\), obtaining an indiscernible sequence
\((c_i:i<\omega)\).  For each \(n\) and each \(I\subseteq\{1,\ldots,n\}\),
the formula
\(\exists y(\bigwedge_{k\in I}\varphi(z_k;y)\wedge\bigwedge_{k\notin I}\neg\varphi(z_k;y))\)
holds of every increasing \(n\)-tuple of the \(a_i\) by (3), hence of every
increasing \(n\)-tuple of the \(c_i\).  So for finite disjoint
\(I,J\subseteq\omega\) the set \(\{\varphi(c_i;y):i\in I\}\cup
\{\neg\varphi(c_j;y):j\in J\}\) is realized, and by (Sat) the set
\(\{\varphi(c_i;y):i\text{ even}\}\cup\{\neg\varphi(c_i;y):i\text{ odd}\}\) is
realized by some \(b\).

\emph{(5)\(\Rightarrow\)(2).}  The \(a_i\) are pairwise distinct: if
\(a_i=a_j\) with \(i<j\), indiscernibility applied to the formula
\(z_1=z_2\) gives \(a_0=a_1\), contradicting \(\varphi(a_0;b)\wedge
\neg\varphi(a_1;b)\).  Let \(n<\omega\) and \(J\subseteq n\); put
\(s(i)=2i\) for \(i\in J\) and \(s(i)=2i+1\) for \(i\notin J\).  Then \(s\) is
strictly increasing and \(s(i)\) is even iff \(i\in J\), so \(b\) witnesses
\(\exists y(\bigwedge_{i\in J}\varphi(a_{s(i)};y)\wedge
\bigwedge_{i\in n\setminus J}\neg\varphi(a_{s(i)};y))\).  By indiscernibility
the same formula holds of \((a_0,\ldots,a_{n-1})\).  So
\(\{a_0,\ldots,a_{n-1}\}\) is shattered.
\end{proof}
\end{annexitem}

\begin{annexitem}{annex:p03-w03-p03-035-sauer-shelah}{thm:p03-w03-p03-035-sauer-shelah}{Exponential or polynomial type counts}
\begin{proof}
\emph{A counting lemma.}  For a finite set \(Y\) and a family \(\mathcal S\)
of subsets of \(Y\), let \(\operatorname{Sh}(\mathcal S)\) be the set of
\(Z\subseteq Y\) shattered by \(\mathcal S\), that is, such that
\(\{S\cap Z:S\in\mathcal S\}\) is the whole power set of \(Z\).  We show
\(\abs{\operatorname{Sh}(\mathcal S)}\geq\abs{\mathcal S}\) by induction on
\(\abs Y\).  If \(Y=\varnothing\), then \(\mathcal S\subseteq\{\varnothing\}\)
and \(\varnothing\) is shattered iff \(\mathcal S\neq\varnothing\).  Otherwise
fix \(y\in Y\), let \(Y'=Y\setminus\{y\}\), and put
\[
  \mathcal S_0=\{S\setminus\{y\}:S\in\mathcal S\},\qquad
  \mathcal S_1=\{S'\subseteq Y':S'\in\mathcal S\text{ and }
  S'\cup\{y\}\in\mathcal S\}.
\]
Each \(S'\in\mathcal S_0\) comes from one member of \(\mathcal S\) if exactly
one of \(S',S'\cup\{y\}\) is in \(\mathcal S\), and from two if
\(S'\in\mathcal S_1\); so \(\abs{\mathcal S}=\abs{\mathcal S_0}+
\abs{\mathcal S_1}\).  A set \(Z\subseteq Y'\) shattered by \(\mathcal S_0\)
is shattered by \(\mathcal S\), since traces on \(Z\) do not see \(y\).  If
\(Z\subseteq Y'\) is shattered by \(\mathcal S_1\), then \(Z\cup\{y\}\) is
shattered by \(\mathcal S\): for \(Z_0\subseteq Z\) pick \(S'\in\mathcal S_1\)
with \(S'\cap Z=Z_0\); then \(S'\) and \(S'\cup\{y\}\) lie in \(\mathcal S\)
and cut out \(Z_0\) and \(Z_0\cup\{y\}\).  The sets of the first kind omit
\(y\) and those of the second contain it, so by the induction hypothesis on
\(Y'\),
\(\abs{\operatorname{Sh}(\mathcal S)}\geq\abs{\mathcal S_0}+\abs{\mathcal S_1}
=\abs{\mathcal S}\).

\emph{(1).}  Let \(B\) be a set of \(n\) distinct \(\abs y\)-tuples and
\(\mathcal S=\{B_p:p\in S_\varphi(B)\}\), with \(B_p\) as in the proof of
\cref{thm:p03-w03-p03-035-ip-characterization}; there it was shown that
\(\abs{\mathcal S}=\abs{S_\varphi(B)}\) and that \(\mathcal S\) is the family of
traces \(\varphi^{\mathrm{opp}}(\mathfrak C;a)\cap B\).  So a subset of \(B\)
shattered by \(\mathcal S\) is shattered by \(\varphi^{\mathrm{opp}}\) and has
at most \(d\) members.  By the counting lemma,
\(\abs{S_\varphi(B)}\leq\abs{\operatorname{Sh}(\mathcal S)}\leq
\sum_{i\leq d}\binom ni\).  For the last inequality, if \(d=0\) both sides are
\(1\); if \(d\geq1\), listing a subset of \(\{1,\ldots,n\}\) with at most \(d\)
elements in increasing order and padding with \(0\) to length \(d\) is an
injection into \(\{0,\ldots,n\}^d\).

\emph{(2).}  If \(\varphi\) is NIP, then \(\varphi^{\mathrm{opp}}\) is NIP
(\cref{thm:p03-w03-p03-035-ip-characterization}, (1)\(\Leftrightarrow\)(4)),
so by (1)\(\Leftrightarrow\)(2) of that theorem some \(n\) has no shattered
set of \(n\) tuples for \(\varphi^{\mathrm{opp}}\); since subsets of
shattered sets are shattered, the VC-dimension \(d\) of
\(\varphi^{\mathrm{opp}}\) is finite and (1) applies.  If
\(\pi_\varphi(n)\leq(n+1)^d\) for all \(n\), then \(\pi_\varphi(n)<2^n\) for
large \(n\).  If \(\pi_\varphi(n)<2^n\) for some \(n\), item~(6) of
\cref{thm:p03-w03-p03-035-ip-characterization} fails, so \(\varphi\) is NIP.
\end{proof}
\end{annexitem}

\begin{annexitem}{annex:p03-w03-p03-035-alternation}{thm:p03-w03-p03-035-alternation}{Bounded alternation on indiscernible sequences}
\begin{proof}
Say that \(i_0<\cdots<i_N\) \emph{alternate for \(b\)} when
\(\varphi(a_{i_k};b)\leftrightarrow\neg\varphi(a_{i_{k+1}};b)\) for all
\(k<N\).  Suppose that for every \(N\) there are an indiscernible sequence
\((a^N_i:i\in I_N)\), a tuple \(b_N\) and indices \(i_0<\cdots<i_N\)
alternating for \(b_N\).  In variables \((x_k:k<\omega)\) and \(y\) let
\(\Sigma\) consist of the formulas
\(\psi(x_{k_1},\ldots,x_{k_n})\leftrightarrow\psi(x_{l_1},\ldots,x_{l_n})\) for
\(\psi\in\Lang\) and increasing \(k_1<\cdots<k_n\), \(l_1<\cdots<l_n\), and the
formulas \(\varphi(x_k;y)\leftrightarrow\neg\varphi(x_{k+1};y)\) for
\(k<\omega\).  A finite subset mentions only \(x_0,\ldots,x_{M-1}\); it is
realized by \(a^M_{i_0},\ldots,a^M_{i_{M-1}}\) and \(b_M\), because an
increasing subsequence of an indiscernible sequence satisfies every formula
of the first kind.  By (Sat), \(\Sigma\) is realized by some \((c_k)\) and
\(b\).  Then \((c_k:k<\omega)\) is indiscernible and \(\varphi(c_k;b)\)
alternates with \(k\).  If \(\varphi(c_0;b)\) holds, \(\varphi(c_k;b)\) holds iff
\(k\) is even; otherwise the indiscernible sequence \((c_{k+1}:k<\omega)\) has
this property.  By \cref{thm:p03-w03-p03-035-ip-characterization},
(5)\(\Rightarrow\)(1), \(\varphi\) has IP, a contradiction.  So some \(N\)
works.

For the equivalence, suppose \(\varphi\) NIP with bound \(N\), and let
\((a_i:i\in I)\), \(I\neq\varnothing\), and \(b\) be given.  Single indices
alternate trivially, and alternating chains have fewer than \(N+1\) members,
so choose an alternating chain \(i_0<\cdots<i_m\) of maximal length.  Let
\(I_0=\{i\in I:i\geq i_m\}\), a nonempty end segment.  If some \(i\in I_0\)
had \(\varphi(a_i;b)\leftrightarrow\neg\varphi(a_{i_m};b)\), then \(i>i_m\)
and \(i_0<\cdots<i_m<i\) would alternate, contrary to maximality.  So
\(\varphi(a_i;b)^{\varepsilon}\) on \(I_0\), where \(\varepsilon\) is the truth
value at \(i_m\).  Conversely, if \(\varphi\) has IP, item~(5) of
\cref{thm:p03-w03-p03-035-ip-characterization} gives an indiscernible
sequence indexed by \(\omega\) and a \(b\) for which every nonempty end
segment contains indices of both parities, so no \(\varepsilon\) exists.
\end{proof}
\end{annexitem}

\begin{annexitem}{annex:p03-w03-p03-035-boolean}{thm:p03-w03-p03-035-boolean}{Boolean combinations of NIP formulas}
\begin{proof}
A set is shattered by \(\neg\varphi\) iff it is shattered by \(\varphi\), since
\(A_0\) is cut out by \(\neg\varphi(x;b)\) exactly when \(A\setminus A_0\) is cut
out by \(\varphi(x;b)\).  If \(z\) does not occur in \(\varphi\), a set is cut
out by \(\varphi(x;bc)\) iff it is cut out by \(\varphi(x;b)\), so
\(\varphi(x;y)\) and \(\varphi(x;yz)\) shatter the same sets.

For \(\varphi\wedge\psi\) use the criterion of
\cref{thm:p03-w03-p03-035-alternation}.  Given an indiscernible
\((a_i:i\in I)\) with \(I\neq\varnothing\) and \(b\), take nonempty end
segments \(I_\varphi\), \(I_\psi\) on which \(\varphi(a_i;b)\), respectively
\(\psi(a_i;b)\), has constant truth value.  End segments of a linear order are
comparable under inclusion: if \(i\in I_\varphi\setminus I_\psi\) and
\(j\in I_\psi\setminus I_\varphi\), then \(i<j\) (else \(i\geq j\) and
\(i\in I_\psi\)) and likewise \(j<i\).  So \(I_\varphi\cap I_\psi\) is one of
them, nonempty, and \((\varphi\wedge\psi)(a_i;b)\) is constant on it.  Hence
\(\varphi\wedge\psi\) is NIP.

For a Boolean combination of NIP formulas \(\varphi_1(x;y_1),\ldots,
\varphi_r(x;y_r)\), regard each \(\varphi_l\) as a formula in
\((x;y_1\cdots y_r)\), which is NIP by the remark on unused variables, and
induct on the construction of the combination from \(\neg\) and \(\wedge\).
\end{proof}
\end{annexitem}

\begin{annexitem}{annex:p03-w03-p03-035-one-variable}{thm:p03-w03-p03-035-one-variable}{NIP theories are detected by one-variable formulas}
\begin{proof}
Clearly (1) implies (2) and (3).  If (3) holds and \(\abs y=1\), then
\(\varphi^{\mathrm{opp}}(y;x)\) has object variable of length one and is NIP,
so \(\varphi\) is NIP by \cref{thm:p03-w03-p03-035-ip-characterization},
(1)\(\Leftrightarrow\)(4); thus (3) implies (2), and symmetrically (2) implies
(3).  It remains to show that (2) implies (1).  Assume (2) and let
\(\kappa=\abs T^+\), a small regular cardinal.

\emph{Claim.}  Let \((a_i:i<\kappa)\) be an indiscernible sequence of tuples
and \(b\) a single element.  Some \(\alpha<\kappa\) has
\((a_i:\alpha<i<\kappa)\) indiscernible over \(b\).

Suppose not.  For each \(\alpha<\kappa\) there are a formula
\(\psi_\alpha(z_1,\ldots,z_n;w)\in\Lang\) with \(\abs w=1\) and increasing
tuples \(\bar\imath^\alpha,\bar\jmath^\alpha\) of indices above \(\alpha\) with
\(\psi_\alpha(a_{\bar\imath^\alpha};b)\wedge\neg\psi_\alpha(a_{\bar\jmath^\alpha};b)\).
There are at most \(\abs T<\kappa\) formulas and \(\kappa\) is regular, so one
formula \(\psi(z_1,\ldots,z_n;w)\) is \(\psi_\alpha\) for all \(\alpha\) in an
unbounded set \(U\subseteq\kappa\).  Choose by induction on \(l<\omega\) some
\(\alpha_l\in U\) above every index used before stage \(l\), and let
\(\bar k^l\) be \(\bar\imath^{\alpha_l}\) for even \(l\) and
\(\bar\jmath^{\alpha_l}\) for odd \(l\).  Let \(d_l\) be the concatenation
\(a_{k^l_1}\cdots a_{k^l_n}\).  For \(l_1<\cdots<l_r\) the concatenation
\(d_{l_1}\cdots d_{l_r}\) is the subsequence of \((a_i)\) at an increasing
tuple of \(nr\) indices, so \((d_l:l<\omega)\) is indiscernible; and
\(\psi(d_l;b)\) holds iff \(l\) is even.  By
\cref{thm:p03-w03-p03-035-ip-characterization}, (5)\(\Rightarrow\)(1), the
partitioned formula \(\psi(z_1\cdots z_n;w)\) has IP, contrary to (2) since
\(\abs w=1\).

\emph{Main argument.}  Let \(\varphi(x;y)\) with \(y=y_1\cdots y_m\) and
suppose \(\varphi\) has IP.  Item~(5) of
\cref{thm:p03-w03-p03-035-ip-characterization} gives an indiscernible
\((a_i:i<\omega)\) and \(b\) with \(\varphi(a_i;b)\) iff \(i\) is even.  Call
an ordinal \(\lambda+2k\), \(\lambda\) zero or a limit, \emph{even}, and the
others odd.  In variables \((x_\alpha:\alpha<\kappa)\) and \(y\), let
\(\Sigma\) say that \((x_\alpha)\) is indiscernible and that
\(\varphi(x_\alpha;y)\) holds iff \(\alpha\) is even.  A finite subset
mentions indices \(\alpha_1<\cdots<\alpha_r\); choose
\(j_1<\cdots<j_r<\omega\) with \(j_l\) of the same parity as \(\alpha_l\) and
assign \(a_{j_l}\) to \(x_{\alpha_l}\) and \(b\) to \(y\).  By (Sat),
\(\Sigma\) is realized by some \((a'_\alpha:\alpha<\kappa)\) and
\(b'=b'_1\cdots b'_m\).  By the claim some \(\beta_1\) has
\((a'_\alpha:\beta_1<\alpha<\kappa)\) indiscernible over \(b'_1\); then the
sequence \((a'_\alpha b'_1:\beta_1<\alpha<\kappa)\) is indiscernible, because a
formula about an increasing tuple of its terms is a formula over \(b'_1\) about
an increasing tuple of the \(a'_\alpha\).  The tail has order type \(\kappa\),
so the claim applies again with \(b'_2\).  After \(m\) steps,
\((a'_\alpha b'_1\cdots b'_m:\beta_m<\alpha<\kappa)\) is indiscernible, so the
truth value of \(\varphi(a'_\alpha;b')\) is the same for all
\(\alpha>\beta_m\).  But \(\beta_m+1\) and \(\beta_m+2\) have different
parities.  This contradiction shows that \(\varphi\) is NIP.

Finally, if \(T\) is NIP and \(\varphi(x;y,d)\) had IP as a partitioned formula
in \((x;y)\), the witnesses \(b\) would give witnesses \(bd\) for the
parameter-free \(\varphi(x;yz)\), which would have IP.
\end{proof}
\end{annexitem}

\begin{annexitem}{annex:p03-w03-p03-035-boolean-algebra-ip}{thm:p03-w03-p03-035-boolean-algebra-ip}{The order of an infinite Boolean algebra has IP}
\begin{proof}
Only the Boolean-algebra laws are used: \(\wedge\) and \(\vee\) are
associative, commutative, absorptive and distribute over each other,
\(x\wedge 1=x\), \(x\vee 0=x\), \(x\wedge\neg x=0\) and \(x\vee\neg x=1\).
The relation \(x\leq y\) (that is, \(x\wedge y=x\)) is a partial order with
least element \(0\), and \(x\leq y\) implies \(x\wedge z\leq y\wedge z\).  For
\(b\in B\) put \(B{\restriction}b=\{x\in B:x\leq b\}\).

\emph{Splitting.}  Let \(b\in B\) and \(a\in B{\restriction}b\), and put
\(a'=b\wedge\neg a\).  For \(x\in B{\restriction}b\),
\[
 x=x\wedge b=x\wedge\bigl((a\vee\neg a)\wedge b\bigr)
  =x\wedge\bigl((a\wedge b)\vee(\neg a\wedge b)\bigr)
  =(x\wedge a)\vee(x\wedge a'),
\]
using \(a\wedge b=a\).  Hence the map
\(x\mapsto(x\wedge a,\;x\wedge a')\) from \(B{\restriction}b\) to
\(B{\restriction}a\times B{\restriction}a'\) is injective: \(x\) is the join
of the two coordinates of its image.  So if \(B{\restriction}b\) is infinite,
at least one of \(B{\restriction}a\), \(B{\restriction}a'\) is infinite,
since a product of two finite sets is finite.  Moreover
\(a\wedge a'=a\wedge\neg a\wedge b=0\).  If \(a\neq 0\) and \(a\neq b\), then
\(a'\neq 0\): from \(a'=0\) we would get
\(b=(b\wedge a)\vee(b\wedge\neg a)=a\vee 0=a\).

\emph{Item \textup{(1)}.}  Fix \(n<\omega\).  Define \(b_0,\ldots,b_n\) and
\(d_0,\ldots,d_{n-1}\) by recursion, keeping the invariant that
\(B{\restriction}b_k\) is infinite.  Put \(b_0=1\); then
\(B{\restriction}b_0=B\) is infinite.  Given \(b_k\) with
\(B{\restriction}b_k\) infinite, that set has an element \(a\) other than
\(0\) and \(b_k\), because it has at least three elements.  With
\(a'=b_k\wedge\neg a\), splitting shows that \(a\) and \(a'\) are nonzero,
\(a\wedge a'=0\), and one of \(B{\restriction}a\), \(B{\restriction}a'\) is
infinite.  Let \(b_{k+1}\) be one of \(a,a'\) whose set is infinite, and
\(d_k\) the other one.  Then \(d_k\neq 0\), \(d_k\leq b_k\),
\(b_{k+1}\leq b_k\) and \(d_k\wedge b_{k+1}=0\).

Now let \(i<j<n\).  The chain \(b_0\geq b_1\geq\cdots\geq b_n\) gives
\(d_j\leq b_j\leq b_{i+1}\), so
\(d_i\wedge d_j\leq d_i\wedge b_{i+1}=0\), that is \(d_i\wedge d_j=0\).  The
\(d_i\) are pairwise distinct: \(d_i=d_j\) with \(i\neq j\) would give
\(d_i=d_i\wedge d_i=d_i\wedge d_j=0\).

\emph{Item \textup{(2)}.}  Let \(I\subseteq n\) and \(j<n\).  If \(j\in I\),
then \(d_j\leq b_I\), because a join of finitely many elements lies above each
of them.  If \(j\notin I\), distributivity gives
\(d_j\wedge b_I=\bigvee_{i\in I}(d_j\wedge d_i)=0\) (the empty join being
\(0\)), so \(d_j\leq b_I\) would give \(d_j=d_j\wedge b_I=0\), contrary to
\(d_j\neq 0\).

\emph{Item \textup{(3)}.}  Let \(A=\{d_0,\ldots,d_{n-1}\}\), a set of \(n\)
elements by \textup{(1)}.  A subset of \(A\) has the form
\(\{d_i:i\in I\}\) for a unique \(I\subseteq n\), and by \textup{(2)} the
parameter \(b_I\) cuts out exactly that subset:
\(\{d\in A:B\entails d\leq b_I\}=\{d_i:i\in I\}\).  So \(A\) is shattered in
\(B\) by \(x\leq y\).  As \(n\) was arbitrary and \(B\) is a model of
\(\operatorname{Th}(B)\), arbitrarily large finite sets are shattered in a
model of \(\operatorname{Th}(B)\), which is IP in the sense of
\cref{def:p03-nip}.
\end{proof}
\end{annexitem}

\begin{annexitem}{annex:p03-w03-p03-035-shelah-rank-basics}{thm:p03-w03-p03-035-shelah-rank-basics}{Shelah's local rank: well-definedness and finite character}
\begin{proof}
Fix \(\varphi\) and suppress its superscripts.  Simultaneous induction on
\(\alpha\) proves two assertions: membership of a fixed type in
\(S_\alpha(A)\) is unchanged on enlarging the ambient parameter set, and
\(p\subseteq q\in S_\alpha(A)\) implies \(p\in S_\alpha(A)\).  At zero
these say that the same realization still realizes a smaller set of
formulas.  At a limit they follow by intersection.  At a successor, every
finite subset of the smaller type is a finite subset of the larger one.
For ambient-set independence, a splitting witness over \(C\supseteq A\)
can be moved to \(C\cup B\) by the induction hypothesis when the ambient
set is enlarged to \(B\); a witness over a superset of \(B\) is already a
witness over a superset of \(A\).  The required membership at the preceding
stage also transfers by induction.  This proves both assertions.

The classes decrease with \(\alpha\), by their definition and induction.
If a realized type leaves them, its least failure stage cannot be zero or
a limit; it is \(\alpha+1\), and its rank is \(\alpha\).  If it never
leaves them its rank is \(\infty\).  This proves (1) and (2).

For (3), first suppose \(\operatorname{Rank}p=\alpha<\infty\).  Failure
of membership at \(\alpha+1\) supplies a finite \(q\subseteq p\) for which
no splitting witness at level \(\alpha\) exists.  Monotonicity gives
\(\operatorname{Rank}q\geq\alpha\), and the same failure prevents its rank
being at least \(\alpha+1\).  Thus its rank equals \(\alpha\).  If
\(\operatorname{Rank}p=\infty\), the finite subset \(q=\varnothing\) also
has infinite rank by monotonicity.  The assertion is about realized types;
an inconsistent set instead has a finite inconsistent subset by compactness.

For (6), every finite \(q\subseteq p\) has both signed extensions at rank
at least \(\alpha\), by (2).  Put all their parameters into one small
ambient set and use (1).  These are the witnesses in the successor clause,
so \(p\) has rank at least \(\alpha+1\).  In (4), monotonicity applied to
\(q_1,q_2\) gives precisely these two signed extensions of \(p\),
contradicting its rank \(\alpha\).  Two distinct complete extensions in
(5) disagree on an instance of \(\varphi\), so (4) proves their uniqueness.
\end{proof}
\end{annexitem}

\begin{annexitem}{annex:p03-w03-p03-035-shelah-rank-trees}{thm:p03-w03-p03-035-shelah-rank-trees}{Shelah's local rank and finite binary trees}
\begin{proof}
We first prove (2) by induction on \(n\).  At zero both conditions say
that \(p\) is realized.  If \(\operatorname{Rank}p\geq n+1\), every finite
\(q\subseteq p\) has a splitting instance for which both signed
extensions have rank at least \(n\).  By induction the two extensions have
trees of height \(n\); joining these below that instance gives a tree of
height \(n+1\) above \(q\).  Every finite subset of \(\Gamma_{p,n+1}\) is
contained in such a tree requirement for some finite \(q\), so (Sat)
realizes \(\Gamma_{p,n+1}\).  Conversely, the root of a tree of height
\(n+1\) leaves trees of height \(n\) above both signed extensions of
\(p\).  The induction hypothesis and
\cref{thm:p03-w03-p03-035-shelah-rank-basics}\,(6) give rank at least
\(n+1\).

Here is the step that rules out an ordinal rank at or above \(\omega\).
Let \(\mathcal H\) be the class of types having trees of every finite
height.  We show by transfinite induction that every member of
\(\mathcal H\) belongs to every \(S_\alpha\).  The zero and limit steps
are immediate.  For the successor step fix \(p\in\mathcal H\) and a
finite \(q\subseteq p\).  There is a single parameter tuple \(b\) for
which both \(q\cup\{\varphi(x;b)\}\) and
\(q\cup\{\neg\varphi(x;b)\}\) belong to \(\mathcal H\): require, in
variables for \(b\) and for countably many trees, trees of height \(n\)
above each of these two signed types, for every \(n\).  A finite part of
this requirement is realized by taking the root of a tree above \(q\)
whose height exceeds every mentioned height, and using its two subtrees.
Compactness and (Sat) give all the trees with the same \(b\).  The
induction hypothesis puts both signed types in \(S_\alpha\), over an
ambient set containing \(b\).  This works for each finite \(q\), proving
\(p\in S_{\alpha+1}\).  Thus every member of \(\mathcal H\) has rank
\(\infty\).

Consistency of \(T\cup\Gamma(\varphi,m,\omega)\) supplies finite trees
of every height above the empty type, so (1) follows.  Conversely, if the
ranks have no finite upper bound, monotonicity and (2) give trees of every
finite height above the empty type.  Every finite part of the infinite
tree scheme is realized in a sufficiently tall finite tree: extend the
finitely many mentioned branches to one common level, preserving all the
branch choices already specified.  Compactness proves consistency of the
infinite scheme.  By (2), rank at least \(\omega\) means membership in
\(\mathcal H\), which we have proved implies infinite rank.  This proves
(3).  Finally (2) identifies the finite rank thresholds with exactly the
thresholds defining \(\mathrm R_\varphi\); the preceding argument
identifies their infinite values as well.  This proves (4).
\end{proof}
\end{annexitem}

\begin{annexitem}{annex:p03-w03-p03-035-shelah-29}{thm:p03-w03-p03-035-shelah-29}{The two cycles characterizing instability of a formula}
\begin{proof}
The implication (A)\(\Rightarrow\)(B) is immediate.  For
(B)\(\Rightarrow\)(C), take the model and the set \(A\) in (B), and pass
to a sufficiently saturated elementary extension.  The complete
\(\varphi\)-types over \(A\) consistent with the original elementary
diagram are unchanged: each finite subset is realized in the original
model, and each such type is realized in the extension.  In particular
there are still more than \(\lambda\) such types, with
\(\abs A\leq\lambda\).  The argument of
\cref{thm:p03-w03-p03-032-many-to-order}, with \(\bar x\) as its object
tuple, therefore gives the order property.  Its finite order patterns are
sentences of the original language, so passing to the extension has not
changed the conclusion about \(\varphi\) in \(T\).

The two clauses of \cref{thm:p03-w03-p03-035-shelah-tree} give
(C)\(\Rightarrow\)(D)\(\Rightarrow\)(A), closing the first cycle.
For the second, \cref{thm:p03-w03-p03-035-shelah-rank-trees}\,(1) gives
(D)\(\Rightarrow\)(E).  Infinite rank is above every finite rank, so
(E)\(\Rightarrow\)(F); and part~(3) of the same lemma gives
(F)\(\Rightarrow\)(D).  Every condition therefore belongs to the same
equivalence class.
\end{proof}
\end{annexitem}


\begin{annexitem}{annex:p03-w03-p03-035-shelah-tree}{thm:p03-w03-p03-035-shelah-tree}{Shelah's Lemmas~2.2 and 2.3: from order to tree to many types}
\begin{proof}
\emph{(1) Placing the tree in the real line.}  For \(\eta\in{}^{\omega}2\) put
\(\ell(\eta)=\sum_{k<\omega}2(1-\eta(k))3^{-k-1}\), and for a node
\(\nu\in{}^{k}2\) put \(s(\nu)=\sum_{i<k}2(1-\nu(i))3^{-i-1}\) and
\(c(\nu)=s(\nu)+\tfrac32\cdot3^{-k-1}\).  If \(\eta\) extends \(\nu\), then
\(\ell(\eta)-s(\nu)=\sum_{i\geq k}2(1-\eta(i))3^{-i-1}\); this lies in
\([0,3^{-k-1}]\) when \(\eta(k)=1\) and in \([2\cdot3^{-k-1},3^{-k}]\) when
\(\eta(k)=0\), because \(\sum_{i>k}2\cdot3^{-i-1}=3^{-k-1}\).  Hence
\[
  \eta(k)=1\iff\ell(\eta)<c(\eta\restriction k).\tag{\(*\)}
\]
The map \(\ell\) is injective (two sequences first differing at \(k\) are
separated as above), every \(\ell(\eta)\) with \(\eta\) extending \(\nu\) lies
in \(I_\nu=[s(\nu),s(\nu)+3^{-k}]\) outside the open middle third, which
contains \(c(\nu)\), and the intervals \(I_\nu\) of incomparable nodes are
disjoint.  So all the reals \(\ell(\eta)\) and \(c(\nu)\) are pairwise
distinct.

\emph{(1) Compactness.}  By the remark after
\cref{def:p03-w03-p03-035-shelah-defs}, \(\varphi\) has the order property of
\cref{def:p03-stability}; as this is expressed by sentences true in a model
of \(T\), there are \((a_i:i<\omega)\), \((b_j:j<\omega)\) in \(\mathfrak C\)
with \(\varphi(a_i;b_j)\) iff \(i<j\).  A finite subset \(\Gamma_0\) of
\(\Gamma(\varphi,m,\omega)\) mentions finitely many leaves \(\eta\) and nodes
\(\nu\).  List the corresponding reals \(\ell(\eta)\), \(c(\nu)\) in
increasing order and let \(t(\cdot)\) be the position in this list; assign
\(a_{t(\ell(\eta))}\) to \(\bar x_\eta\) and \(b_{t(c(\nu))}\) to
\(\bar y_\nu\).  For a formula
\(\varphi(\bar x_\eta;\bar y_{\eta\restriction k})^{\eta(k)}\) of
\(\Gamma_0\), \(\varphi(a_{t(\ell(\eta))};b_{t(c(\eta\restriction k))})\)
holds iff \(t(\ell(\eta))<t(c(\eta\restriction k))\) iff
\(\ell(\eta)<c(\eta\restriction k)\) iff \(\eta(k)=1\), by \((*)\).  So
\(\Gamma_0\) is satisfied in \(\mathfrak C\), and by compactness
\(T\cup\Gamma(\varphi,m,\omega)\) is consistent.

\emph{(2) A longer tree.}  Let \(\mu\) be the least cardinal with
\(2^\mu>\lambda\); then \(\mu\leq\lambda\), since \(2^\lambda>\lambda\),
and \(2^{\abs\kappa}\leq\lambda\) for every ordinal \(\kappa<\mu\).  We show that \(T\cup\Gamma(\varphi,m,\mu)\) is
consistent.  A finite subset mentions finitely many leaves
\(\eta\in{}^{\mu}2\) and finitely many levels.  Let \(K=\{k_0<\cdots<
k_{r-1}\}\) be a finite set of ordinals below \(\mu\) containing these levels
and, for each pair of distinct mentioned leaves, the least level at which
they differ.  Send a
leaf \(\eta\) to \(\eta^*\in{}^{\omega}2\) with \(\eta^*(i)=\eta(k_i)\) for
\(i<r\) and \(\eta^*(i)=0\) for \(i\geq r\), and the node
\(\eta\restriction k_i\) to \(\eta^*\restriction i\).  This is well defined
and injective on the mentioned nodes: if \(\eta\restriction k_i=
\eta'\restriction k_i\) then \(\eta(k_j)=\eta'(k_j)\) for \(j<i\); if not, the
least level where \(\eta,\eta'\) differ is some \(k_j<k_i\) in \(K\), and
\(\eta^*(j)\neq\eta'^*(j)\).  It sends \(\varphi(\bar x_\eta;\bar
y_{\eta\restriction k_i})^{\eta(k_i)}\) to
\(\varphi(\bar x_{\eta^*};\bar y_{\eta^*\restriction i})^{\eta^*(i)}\), a
member of \(\Gamma(\varphi,m,\omega)\).  So each finite subset is a renaming of
variables in a finite subset of \(\Gamma(\varphi,m,\omega)\), and
compactness gives a model \(\mathcal M\) of \(T\cup\Gamma(\varphi,m,\mu)\),
with \(\bar a_\nu\) interpreting \(\bar y_\nu\) and \(\bar c_\eta\)
interpreting \(\bar x_\eta\).

\emph{(2) Counting.}  Let \(A\) be the set of entries of the \(\bar a_\nu\)
with \(\nu\in{}^{\kappa}2\), \(\kappa<\mu\).  There are
\(\sum_{\kappa<\mu}2^{\abs\kappa}\leq\mu\cdot\lambda=\lambda\) such nodes, each
contributing finitely many entries, so \(\abs A\leq\lambda\).  For
\(\eta\in{}^{\mu}2\) let \(p_\eta\) be the complete \(\varphi\)-\(m\)-type of
\(\bar c_\eta\) over \(A\).  If \(\eta\neq\eta'\) first differ at \(k<\mu\),
then \(\eta\restriction k=\eta'\restriction k\) and \(p_\eta\) contains
\(\varphi(\bar x;\bar a_{\eta\restriction k})^{\eta(k)}\) while \(p_{\eta'}\)
contains its negation, so \(p_\eta\neq p_{\eta'}\).  Hence
\(\abs{S^m_\varphi(A)}\geq2^\mu>\lambda\).
\end{proof}
\end{annexitem}

\begin{annexitem}{annex:p03-w03-p03-035-dividing-basics}{thm:p03-w03-p03-035-dividing-basics}{Dividing: finite inconsistency and smaller bases}
\begin{proof}
By \cref{thm:p03-w03-p03-031-saturation}~(3) the witnesses may be taken in
\(\mathfrak C\).
(1) Let \((b_i:i<\omega)\) be \(A\)-indiscernible with \(b_0=b\) and
\(\{\varphi(x;b_i)\}\) inconsistent.  By compactness some finite subset is
inconsistent, say for \(i_1<\cdots<i_k\); so
\(\neg\exists x\bigwedge_{l\leq k}\varphi(x;z_l)\) holds of
\((b_{i_1},\ldots,b_{i_k})\), hence by indiscernibility of every increasing
\(k\)-tuple, and every \(k\)-element subset is inconsistent.
(2) A \(C\)-indiscernible sequence is \(A\)-indiscernible, since every
\(\Lang(A)\)-formula is an \(\Lang(C)\)-formula; the same sequence witnesses
dividing over \(A\).
(3) Take \(m=1\) and \(\varphi_1=\varphi\) in
\cref{def:p03-dividing-forking}~(2).  A complete type containing a formula that
divides over \(A\) therefore contains a formula that forks over \(A\).
\end{proof}
\end{annexitem}

\begin{annexitem}{annex:p03-w03-p03-035-sop-not-simple}{thm:p03-w03-p03-035-sop-not-simple}{The strict order property excludes simplicity}
\begin{proof}
Let \(\psi(x;y)\) and \((b_i:i<\omega)\) witness SOP.  Put
\(P(y_1,y_2)\defeq\forall x(\psi(x;y_1)\rightarrow\psi(x;y_2))\) and
\(\rho(y_1,y_2)\defeq P(y_1,y_2)\wedge\neg P(y_2,y_1)\).  \(P\) is reflexive
and transitive, so \(\rho\) is irreflexive, and it is transitive: from
\(\rho(a,b)\) and \(\rho(b,c)\) we get \(P(a,c)\), and \(P(c,a)\) would give
\(P(b,a)\) from \(P(b,c)\).  For \(i<j\),
\(\psi(\mathfrak C;b_i)\subsetneq\psi(\mathfrak C;b_{i+1})\subseteq
\psi(\mathfrak C;b_j)\), so \(\rho(b_i,b_j)\).

Let \(\kappa=\abs T^+\).  Let \(I\) be the linear order whose elements are
\(p_\alpha,q_\alpha\) (\(\alpha<\kappa\)) and \(r_{\alpha,j},s_{\alpha,j}\)
(\(\alpha<\kappa\), \(1\leq j<\omega\)), ordered as follows: the \(p_\alpha\)
increase with \(\alpha\) and lie below all other elements; for each \(\alpha\)
the \emph{block} \(q_\alpha<r_{\alpha,1}<s_{\alpha,1}<r_{\alpha,2}<
s_{\alpha,2}<\cdots\) is ordered as written; and for \(\beta<\alpha\) the whole
block of \(\beta\) lies above the whole block of \(\alpha\).  \(I\) is small.
By \cref{thm:p03-w03-p03-035-em-extraction}, applied to \((b_i)\) with
\(A=\varnothing\) and \(X=I\), there is an indiscernible
\((e_t:t\in I)\) with \(\rho(e_s,e_t)\) whenever \(s<t\).

Let \(u\) be a tuple of variables of length \(\abs y\) and
\(\chi(u;v_1v_2)\defeq\rho(v_1,u)\wedge\rho(u,v_2)\).  Put
\(c_\alpha=e_{p_\alpha}e_{q_\alpha}\), \(c^0_\alpha=c_\alpha\),
\(c^j_\alpha=e_{r_{\alpha,j}}e_{s_{\alpha,j}}\) for \(j\geq1\), and let \(B\) and
\(B_\alpha\) be the sets of entries of the \(c_\beta\) for \(\beta<\kappa\),
respectively \(\beta<\alpha\).

\emph{Consistency.}  For \(\alpha_1<\cdots<\alpha_m\) and \(\beta\) above
them, \(p_{\alpha_l}<p_\beta<q_{\alpha_l}\) in \(I\), so \(e_{p_\beta}\) satisfies
every \(\chi(u;c_{\alpha_l})\).  By (Sat) some \(a\) satisfies all
\(\chi(u;c_\alpha)\); let \(p=\operatorname{tp}(a/B)\).

\emph{Indiscernibility over \(B_\alpha\).}  Relative to the indices
\(p_\beta,q_\beta\) (\(\beta<\alpha\)) of \(B_\alpha\), every index of every
\(c^j_\alpha\) lies above all \(p_\beta\) and below all \(q_\beta\), and among
themselves they are ordered
\(p_\alpha<q_\alpha<r_{\alpha,1}<s_{\alpha,1}<r_{\alpha,2}<\cdots\).  Hence for
\(j_1<\cdots<j_n\), finitely many elements \(\bar d\) of \(B_\alpha\), and any
\(j'_1<\cdots<j'_n\), the index tuples of \(\bar d\,c^{j_1}_\alpha\cdots
c^{j_n}_\alpha\) and \(\bar d\,c^{j'_1}_\alpha\cdots c^{j'_n}_\alpha\) have the
same order type, and by indiscernibility of \((e_t)\) these tuples have the
same type.  So \((c^j_\alpha:j<\omega)\) is \(B_\alpha\)-indiscernible, and its
first member is \(c_\alpha\).

\emph{Inconsistency.}  For \(j<j'\), the second index of \(c^j_\alpha\) is
below the first index of \(c^{j'}_\alpha\).  If \(u\) satisfied
\(\chi(u;c^j_\alpha)\wedge\chi(u;c^{j'}_\alpha)\), transitivity of \(\rho\)
along \(u\), the second entry of \(c^j_\alpha\), the first entry of
\(c^{j'}_\alpha\), and back to \(u\) would give \(\rho(u,u)\).  So
\(\{\chi(u;c^j_\alpha):j<\omega\}\) is \(2\)-inconsistent, and
\(\chi(u;c_\alpha)\) divides over \(B_\alpha\).

\emph{Failure of local character.}  Let \(A\subseteq B\) with
\(\abs A\leq\abs T\).  Each element of \(A\) lies in some \(c_\beta\); as
\(\kappa\) is regular and \(\abs A<\kappa\), these \(\beta\) are bounded by some
\(\alpha<\kappa\), and \(A\subseteq B_\alpha\).  By
\cref{thm:p03-w03-p03-035-dividing-basics}, \(\chi(u;c_\alpha)\) divides and
forks over \(A\), so \(p\) forks over \(A\).  Since \(\abs T\) is infinite,
\(\abs T+\abs u=\abs T\), and \(p\) violates the condition of
\cref{def:p03-simple-theory} in the model \(\mathfrak C\).
\end{proof}
\end{annexitem}

\begin{annexitem}{annex:p03-w03-p03-035-stable-nip}{thm:p03-w03-p03-035-stable-nip}{Stable theories are NIP}
\begin{proof}
Let \(\varphi(x;y)\) have IP and take \((a_i:i<\omega)\) as in
\cref{thm:p03-w03-p03-035-ip-characterization}~(3).  For \(j<\omega\) the set
\(\{\varphi(a_i;y):i<j\}\cup\{\neg\varphi(a_i;y):i\geq j\}\) is finitely
realized by (3) and has countably many parameters, so it is realized by some
\(b_j\) by (Sat).  Then \(\mathfrak C\entails\varphi(a_i;b_j)\) iff \(i<j\),
which is the order property of \cref{def:p03-stability}.  So a theory with an
IP formula is unstable, and a stable theory is NIP.
\end{proof}
\end{annexitem}

\begin{annexitem}{annex:p03-w03-p03-035-ip-or-sop}{thm:p03-w03-p03-035-ip-or-sop}{Instability splits into independence and strict order}
\begin{proof}
Let \(T\) be unstable and \(\varphi(x;y)\) have the order property, witnessed
by \((a_i)\), \((b_j)\) with \(\varphi(a_i;b_j)\) iff \(i<j\); the witnesses can
be moved into \(\mathfrak C\) by (Sat).  If \(\varphi\) has IP we are done, so
assume \(\varphi\) is NIP.

\emph{Step 1: an indiscernible ladder.}  Apply
\cref{thm:p03-w03-p03-035-em-extraction} to the pairs \((a_ib_i:i<\omega)\)
with \(A=\varnothing\) and \(X=\mathbb Q\).  The Ehrenfeucht--Mostowski type
of the pairs contains \(\varphi(z^x_1;z^y_2)\), \(\neg\varphi(z^x_2;z^y_1)\)
and \(\neg\varphi(z^x_1;z^y_1)\), writing \(z_k=z^x_kz^y_k\).  So we get an
indiscernible \((a'_tb'_t:t\in\mathbb Q)\) with \(\varphi(a'_s;b'_t)\) iff
\(s<t\); in particular \((a'_t:t\in\mathbb Q)\) is indiscernible, and the
\(a'_t\) are distinct, since for \(s<s'\) the tuple \(b'_{s'}\) separates
\(a'_s\) from \(a'_{s'}\).

\emph{Step 2: a forbidden pattern.}  As \(\varphi\) is NIP, by
\cref{thm:p03-w03-p03-035-ip-characterization}, (1)\(\Leftrightarrow\)(2),
there is \(n\) such that no set of \(n\) distinct tuples is shattered.  Fix
\(t_0<\cdots<t_{n-1}\) in \(\mathbb Q\); then some
\(\eta:n\to\{0,1\}\) makes \(\{\varphi(a'_{t_i};y)^{\eta(i)}:i<n\}\)
unrealized.  For each \(N\leq n\), the pattern \(1^N0^{n-N}\) (value \(1\)
below \(N\), \(0\) from \(N\) on) is realized, by \(b'_t\) for any \(t\) with
\(t_{N-1}<t<t_N\), the missing bounds being dropped.

\emph{Step 3: a switching position.}  Starting from \(\eta\), repeatedly
replace an adjacent pair of values \(0,1\) at positions \(i,i+1\) by
\(1,0\).  Each replacement lowers the number of pairs \(i<i'\) with
\(\eta(i)=0\), \(\eta(i')=1\) by one, so the process stops, at the pattern
\(1^N0^{n-N}\) with \(N\) the number of \(1\)'s of \(\eta\).  The first pattern
is unrealized and the last realized, so for some consecutive patterns
\(\eta^s,\eta^{s+1}\), switching at positions \(i_0,i_0+1\), \(\eta^s\) is
unrealized and \(\eta^{s+1}\) is realized.  Let
\(\theta_0(y)=\bigwedge_{i\notin\{i_0,i_0+1\}}\varphi(a'_{t_i};y)^{\eta^s(i)}\).
Then
\[
  \theta_0(y)\wedge\neg\varphi(a'_{t_{i_0}};y)\wedge\varphi(a'_{t_{i_0+1}};y)
  \text{ is unrealized,}\qquad
  \theta_0(y)\wedge\varphi(a'_{t_{i_0}};y)\wedge\neg\varphi(a'_{t_{i_0+1}};y)
  \text{ is realized.}
\]

\emph{Step 4: a strictly increasing chain.}  Let \(J\) be the open interval of
\(\mathbb Q\) between \(t_{i_0-1}\) and \(t_{i_0+2}\), a missing bound being
dropped.  For \(s<s'\) in \(J\), the tuple of the parameters of \(\theta_0\)
with \(a'_s,a'_{s'}\) in positions \(i_0,i_0+1\) has the same order type of
indices as the original, so by indiscernibility the two displayed statements
hold with \(a'_s,a'_{s'}\) in place of \(a'_{t_{i_0}},a'_{t_{i_0+1}}\).  Let
\(\zeta(y;x)=\theta_0(y)\wedge\varphi(x;y)\).  For \(s<s'\) in \(J\): if
\(y\in\zeta(\mathfrak C;a'_{s'})\) then \(\varphi(a'_s;y)\), since otherwise
the unrealized combination would be realized; so
\(\zeta(\mathfrak C;a'_{s'})\subseteq\zeta(\mathfrak C;a'_s)\), and a
realization of the realized combination lies in
\(\zeta(\mathfrak C;a'_s)\setminus\zeta(\mathfrak C;a'_{s'})\).  Choose a
strictly decreasing sequence \((s_k:k<\omega)\) in \(J\).  Then
\(\zeta(\mathfrak C;a'_{s_k})\subsetneq\zeta(\mathfrak C;a'_{s_{k+1}})\) for
all \(k\).

\emph{Step 5: removing parameters.}  The parameters of \(\theta_0\) are the
\(n-2\) tuples \(a'_{t_i}\), \(i\notin\{i_0,i_0+1\}\); let \(\theta(y;z)\) be
\(\theta_0\) with these replaced by variables \(z\) and \(\bar a^*\) their
concatenation.  The partitioned formula
\(\theta(y;z)\wedge\varphi(x;y)\), with object variables \(y\) and parameter
variables \(xz\), has the chain of instances given by the parameters
\(a'_{s_k}\bar a^*\), so it has SOP.
\end{proof}
\end{annexitem}

\begin{annexitem}{annex:p03-w03-p03-035-simple-nip-stable}{thm:p03-w03-p03-035-simple-nip-stable}{Simple and NIP theories are stable}
\begin{proof}
Let \(T\) be simple and NIP, and suppose \(T\) unstable.  By
\cref{thm:p03-w03-p03-035-ip-or-sop} some formula has IP or some formula has
SOP.  The first contradicts NIP; the second contradicts simplicity by
\cref{thm:p03-w03-p03-035-sop-not-simple}.  The converse sentence is
\cref{thm:p03-w03-p03-035-stable-nip}.
\end{proof}
\end{annexitem}

\begin{annexitem}{annex:p03-w03-p03-035-tree-order}{thm:p03-w03-p03-035-tree-order}{The tree property gives the order property}
\begin{proof}
We define \(\nu\in\omega^\omega\) and tuples \(e_l\) (\(l<\omega\)) by
induction on \(l\).  Suppose \(\nu\restriction l\) and \(e_0,\ldots,e_{l-1}\)
are defined.  For each \(m<l\), at most \(k-1\) indices \(i\) have
\(\varphi(e_m;b_{(\nu\restriction l)\frown\langle i\rangle})\), because these
siblings are \(k\)-inconsistent.  Choose \(\nu(l)\) outside these finitely
many indices, for all \(m<l\), and choose any \(i_l\neq\nu(l)\).  Let
\(\eta_l\) be the branch \((\nu\restriction l)\frown\langle
i_l\rangle\frown\langle0,0,\ldots\rangle\) and let \(e_l\) realize
\(\{\varphi(x;b_{\eta_l\restriction j}):1\leq j<\omega\}\).  Since
\(\eta_l\restriction j=\nu\restriction j\) for \(j\leq l\), \(e_l\) satisfies
\(\varphi(x;b_{\nu\restriction j})\) for \(1\leq j\leq l\).

Let \(s_l=b_{\nu\restriction(l+1)}\).  If \(l<m\), then \(l+1\leq m\) and
\(\varphi(e_m;s_l)\) holds.  If \(m<l\), then
\(s_l=b_{(\nu\restriction l)\frown\langle\nu(l)\rangle}\), and \(\nu(l)\) was
chosen so that \(\neg\varphi(e_m;s_l)\).  Put \(a_i=e_{2i+2}\) and
\(d_j=s_{2j+1}\).  The indices \(2i+2\) and \(2j+1\) are never equal, and
\(2j+1<2i+2\) iff \(j\leq i\).  So \(\neg\varphi(a_i;d_j)\) holds iff \(i<j\):
\(\neg\varphi(x;y)\) has the order property of \cref{def:p03-stability}.
\end{proof}
\end{annexitem}

\begin{annexitem}{annex:p03-w03-p03-035-stable-dividing}{thm:p03-w03-p03-035-stable-dividing}{Stable theories have local character of dividing}
\begin{proof}
Suppose \(p\) is a complete type over a small \(B\), in finitely many
variables \(x\), that divides over every \(A\subseteq B\) with
\(\abs A\leq\abs T\).  We show that \(T\) is unstable.

\emph{A dividing chain.}  Let \(\kappa=\abs T^+\).  For \(\alpha<\kappa\)
choose by induction a formula \(\varphi_\alpha(x;b_\alpha)\in p\), with
\(b_\alpha\) from \(B\), dividing over the set \(A_\alpha\) of entries of the
\(b_\beta\), \(\beta<\alpha\); this is possible since
\(\abs{A_\alpha}\leq\abs T\).  By
\cref{thm:p03-w03-p03-035-dividing-basics}~(1) there is \(k_\alpha\) and an
\(A_\alpha\)-indiscernible sequence starting at \(b_\alpha\) whose instances
are \(k_\alpha\)-inconsistent.  There are at most \(\abs T\) pairs
\((\varphi_\alpha,k_\alpha)\) and \(\kappa\) is regular, so there are
\(\alpha_0<\alpha_1<\cdots\) with the same pair \((\varphi,k)\).  Put
\(d_n=b_{\alpha_n}\) and let \(D_n\) be the set of entries of
\(d_0,\ldots,d_{n-1}\).  As \(D_n\subseteq A_{\alpha_n}\), the sequence for
\(b_{\alpha_n}\) is \(D_n\)-indiscernible, so there are \(d_n^i\)
(\(i<\omega\)) with \(d_n^0=d_n\), \(d_n^i\equiv_{D_n}d_n\), and
\(\{\varphi(x;d_n^i):i<\omega\}\) \(k\)-inconsistent.  The set
\(\{\varphi(x;d_n):n<\omega\}\subseteq p\) is realized.  Also \(k\geq2\), since
\(\varphi(x;d_0)\in p\) is realized.

\emph{Finite trees.}  Fix \(n\).  We show by downward induction on
\(j=n,n-1,\ldots,0\) the statement \(S_j\): there is a family
\((e_\eta:\eta\in\omega^{\leq n-j},\ \eta\neq\varnothing)\) such that
\begin{enumerate}
  \item for every \(\eta\in\omega^{n-j}\) the set
        \(\{\varphi(x;d_m):m<j\}\cup\{\varphi(x;e_{\eta\restriction l}):
        1\leq l\leq n-j\}\) is realized; and
  \item for every \(\eta\in\omega^{<n-j}\) the set
        \(\{\varphi(x;e_{\eta\frown\langle i\rangle}):i<\omega\}\) is
        \(k\)-inconsistent.
\end{enumerate}
\(S_n\) holds with the empty family, since \(\{\varphi(x;d_m):m<n\}\) is
realized.  Assume \(S_{j+1}\) with family \(F=(e_\eta)\).  For each \(i\),
\(d_j\equiv_{D_j}d^i_j\), so by (Copy) there is \(F^i\) with
\(d_jF\equiv_{D_j}d^i_jF^i\).  Define \(e'_{\langle i\rangle}=d^i_j\) and
\(e'_{\langle i\rangle\frown\eta}=F^i_\eta\).  For a sequence
\(\langle i\rangle\frown\eta\) of length \(n-j\), item~(1) is the truth of
\(\exists x(\bigwedge_{m<j}\varphi(x;d_m)\wedge\varphi(x;d^i_j)\wedge
\bigwedge_l\varphi(x;F^i_{\eta\restriction l}))\), a formula over \(D_j\) about
\(d^i_jF^i\) which holds of \(d_jF\) by \(S_{j+1}\).  Item~(2) holds at the
root because the \(\varphi(x;d^i_j)\) are \(k\)-inconsistent, and inside each
\(F^i\) because \(k\)-inconsistency of finitely many instances is expressed by
a formula without parameters that holds in \(F\).  So \(S_j\) holds.

\emph{An infinite tree.}  In variables \((z_\eta:\eta\in\omega^{<\omega},
\eta\neq\varnothing)\) let \(\Sigma\) consist of
\(\exists x\bigwedge_{1\leq l\leq L}\varphi(x;z_{\eta\restriction l})\) for
\(\eta\in\omega^L\), \(L\geq1\), and
\(\neg\exists x\bigwedge_{i\in w}\varphi(x;z_{\eta\frown\langle i\rangle})\) for
\(\eta\in\omega^{<\omega}\) and \(w\subseteq\omega\) with \(k\) elements.  A
finite subset mentions only nodes of length at most some \(n\), and the family
given by \(S_0\) for that \(n\) realizes it: paths of length \(L\leq n\) are
initial parts of paths of length \(n\).  By (Sat), \(\Sigma\) is realized by
some \((b_\eta)\).  For \(\nu\in\omega^\omega\) the set
\(\{\varphi(x;b_{\nu\restriction l}):l\geq1\}\) is finitely realized, hence
realized by (Sat).  So \(\varphi\) and \(k\) satisfy the hypothesis of
\cref{thm:p03-w03-p03-035-tree-order}, and \(T\) is unstable.
\end{proof}
\end{annexitem}

\begin{annexitem}{annex:p03-w03-p03-035-drank-calculus}{thm:p03-w03-p03-035-drank-calculus}{Calculus of the local dividing rank}
\begin{proof}
Each item is proved by induction on \(n\), simultaneously for all formulas
of the stated shape.

(1) For \(n=0\): a witness \((d_l)\) for \(\operatorname D(\theta)\geq1\) has
\(\theta\wedge\delta(x;d_0)\) realized, so \(\theta\) is realized.  If
\(\operatorname D(\theta)\geq n+2\) is witnessed by \((d_l)\), then
\(\operatorname D(\theta\wedge\delta(x;d_l))\geq n+1\), hence \(\geq n\) by
the induction hypothesis, for every \(l\), and \((d_l)\) witnesses
\(\operatorname D(\theta)\geq n+1\).

(2) For \(n=0\), a realization of \(\theta\) realizes \(\theta'\).  If
\((d_l)\) witnesses \(\operatorname D(\theta)\geq n+1\), then
\(\theta\wedge\delta(x;d_l)\) implies \(\theta'\wedge\delta(x;d_l)\) in
\(\mathfrak C\), so the induction hypothesis gives
\(\operatorname D(\theta'\wedge\delta(x;d_l))\geq n\) for every \(l\), and
\((d_l)\) witnesses \(\operatorname D(\theta')\geq n+1\).

(3) By symmetry of \(\equiv\) it suffices to prove one direction.  For
\(n=0\): \(\theta(x;e)\) is realized iff
\(\mathfrak C\entails\exists x\,\theta(x;e)\), and this formula is in the
type of \(e\), which is the type of \(e'\).  Let \((d_l:l<\omega)\) witness
\(\operatorname D(\theta(x;e))\geq n+1\).  The sequence \((d_l)\) is a small
tuple, so by (Copy) with base \(\varnothing\) there are \((d'_l:l<\omega)\)
with \(e\,(d_l)_{l<\omega}\equiv e'\,(d'_l)_{l<\omega}\).  For
\(l_1<\cdots<l_k\) the formula
\(\neg\exists x\bigwedge_{r=1}^{k}\delta(x;y_r)\) holds of
\((d_{l_1},\ldots,d_{l_k})\), hence of \((d'_{l_1},\ldots,d'_{l_k})\); applied
to every arrangement of \(k\) distinct indices, this shows that
\((\delta(x;d'_l):l<\omega)\) is \(k\)-inconsistent.  For each \(l\),
\(ed_l\equiv e'd'_l\), and the induction hypothesis for the formula
\(\theta(x;z)\wedge\delta(x;y)\) gives
\(\operatorname D(\theta(x;e')\wedge\delta(x;d'_l))\geq n\).  So
\(\operatorname D(\theta(x;e'))\geq n+1\).

(4) By induction on \(m\) it suffices to treat \(m=2\), since
\(\theta_1\vee\cdots\vee\theta_m\) is \((\theta_1\vee\cdots\vee
\theta_{m-1})\vee\theta_m\).  For \(n=0\), a realization of
\(\theta_1\vee\theta_2\) realizes one of them.  Let \((d_l)\) witness
\(\operatorname D(\theta_1\vee\theta_2)\geq n+1\).  For each \(l\), the formula
\((\theta_1\vee\theta_2)\wedge\delta(x;d_l)\) is equivalent in
\(\mathfrak C\) to
\((\theta_1\wedge\delta(x;d_l))\vee(\theta_2\wedge\delta(x;d_l))\), so by (2)
and the induction hypothesis there is \(j_l\in\{1,2\}\) with
\(\operatorname D(\theta_{j_l}\wedge\delta(x;d_l))\geq n\).  Some
\(j\in\{1,2\}\) equals \(j_l\) for infinitely many \(l\); list these \(l\) in
increasing order as \(l(0)<l(1)<\cdots\).  The subfamily
\((\delta(x;d_{l(r)}):r<\omega)\) is \(k\)-inconsistent, because distinct
\(r\) give distinct \(l(r)\).  So \((d_{l(r)}:r<\omega)\) witnesses
\(\operatorname D(\theta_j)\geq n+1\).

(5) First, if \(\operatorname D(\theta)\geq n\), there is a family
\((e_\eta:\eta\in\omega^{\leq n},\ \eta\neq\varnothing)\) of \(\abs y\)-tuples
such that (i) for every \(\eta\in\omega^{n}\) the formula
\(\theta\wedge\bigwedge_{1\leq l\leq n}\delta(x;e_{\eta\restriction l})\) is
realized, and (ii) for every \(\eta\in\omega^{<n}\) the indexed family
\((\delta(x;e_{\eta\frown\langle i\rangle}):i<\omega)\) is
\(k\)-inconsistent.  For \(n=0\) the family is empty and (i) says that
\(\theta\) is realized.  If \((d_i)\) witnesses
\(\operatorname D(\theta)\geq n+1\), the induction hypothesis applied to
\(\theta\wedge\delta(x;d_i)\) gives a family \(F^i\) for each \(i\); put
\(e_{\langle i\rangle}=d_i\) and \(e_{\langle i\rangle\frown\eta}=F^i_\eta\).
Clause (i) for \(\langle i\rangle\frown\eta\) is clause (i) of \(F^i\) for
\(\eta\), clause (ii) at the root is the \(k\)-inconsistency of
\((\delta(x;d_i))\), and clause (ii) below \(\langle i\rangle\) is clause
(ii) of \(F^i\).

Now suppose \(T\) stable and \(\operatorname D(x=x)\geq n\) for every \(n\).
If \(k=1\), a witness \((d_l)\) for \(\operatorname D(x=x)\geq1\) would have
\(\delta(x;d_0)\) realized, since its rank is at least \(0\), and
unrealized, by \(1\)-inconsistency; so \(k\geq2\).  In
variables \((z_\eta:\eta\in\omega^{<\omega},\ \eta\neq\varnothing)\) let
\(\Sigma\) consist of the formulas
\(\exists x\bigwedge_{1\leq l\leq L}\delta(x;z_{\eta\restriction l})\) for
\(\eta\in\omega^{L}\), \(L\geq1\), and
\(\neg\exists x\bigwedge_{r=1}^{k}\delta(x;z_{\eta\frown\langle i_r\rangle})\)
for \(\eta\in\omega^{<\omega}\) and pairwise distinct
\(i_1,\ldots,i_k<\omega\).  A finite subset of \(\Sigma\) mentions only nodes
of length at most some \(n\); the family obtained above for \(x=x\) and that
\(n\) satisfies it, since a node of length \(L\leq n\) is an initial segment
of a node of length \(n\).  The variables are countably many tuples and
there are no parameters, so by (Sat) \(\Sigma\) is realized by some
\((b_\eta)\).  For \(\nu\in\omega^\omega\), every finite part of
\(\{\delta(x;b_{\nu\restriction l}):l\geq1\}\) is realized, so the set is
realized by (Sat).  Thus \(\delta\) and \(k\) satisfy the hypothesis of
\cref{thm:p03-w03-p03-035-tree-order}, and \(T\) is unstable, a
contradiction.  Hence \(\operatorname D(x=x)\) is finite, and
\(\operatorname D(\theta)\leq\operatorname D(x=x)\) by (2).
\end{proof}
\end{annexitem}

\begin{annexitem}{annex:p03-w03-p03-035-stable-simple}{thm:p03-w03-p03-035-stable-simple}{Stable theories are simple}
\begin{proof}
\emph{Finitely many variables.}  Let \(p\) be a complete type over a small
\(B\) in the finite tuple \(x\); it is realized by (Sat), so every formula of
\(p\) has rank at least \(0\).  A pair \((\delta,k)\) consists of a
partitioned formula \(\delta(x;y)\) without parameters, with object
variables \(x\) and some finite parameter tuple \(y\), and an integer
\(k\geq1\); there are \(\abs T\) such pairs.  For each pair, by
\cref{thm:p03-w03-p03-035-drank-calculus}~(5) the values
\(\operatorname D_{\delta,k}(\theta)\) for \(\theta\in p\) are natural
numbers, so we may choose \(\theta_{\delta,k}\in p\) of least
\(\operatorname D_{\delta,k}\)-rank.  Let \(A\) be the set of parameters
occurring in the formulas \(\theta_{\delta,k}\); then \(A\subseteq B\) and
\(\abs A\leq\abs T\).

Suppose some \(\varphi(x;b)\in p\) forks over \(A\).  By
\cref{def:p03-dividing-forking} and \cref{thm:p03-w03-p03-031-saturation}~(3)
there are formulas \(\psi_1(x;c_1),\ldots,\psi_m(x;c_m)\) with parameters in
\(\mathfrak C\), each dividing over \(A\), such that
\(\mathfrak C\entails\forall x(\varphi(x;b)\rightarrow\bigvee_{j}\psi_j(x;c_j))\).
By \cref{thm:p03-w03-p03-035-dividing-basics}~(1), for each \(j\) there are
\(k_j\) and an \(A\)-indiscernible sequence \((c^l_j:l<\omega)\) with
\(c^0_j=c_j\) and \((\psi_j(x;c^l_j):l<\omega)\) \(k_j\)-inconsistent.  Let
\(k=\max_j k_j\); each of these families is \(k\)-inconsistent, since \(k\)
distinct indices include \(k_j\) distinct ones.

\emph{Coding the disjuncts.}  Let \(y_1,\ldots,y_m\) be disjoint tuples of
variables with \(\abs{y_j}=\abs{c_j}\), let \(w,w_1,\ldots,w_m\) be single
variables, \(y=y_1\cdots y_m\,w\,w_1\cdots w_m\), and
\[
  \delta(x;y)\defeq\bigvee_{j=1}^{m}\Bigl(w=w_j\wedge
  \bigwedge_{j'\neq j}w\neq w_{j'}\wedge\psi_j(x;y_j)\Bigr).
\]
Fix pairwise distinct elements \(g_1,\ldots,g_m\) of \(\mathfrak C\), which is
infinite, and an element \(g\).  For \(j\leq m\) and a tuple \(c\) of length
\(\abs{c_j}\), let \(\hat c^{(j)}\) be the \(\abs y\)-tuple with \(c\) in the
slot \(y_j\), \(g\) in every entry of the other slots \(y_{j'}\), \(g_j\) in
the slot \(w\) and \(g_{j'}\) in the slot \(w_{j'}\).  In \(\delta(x;\hat
c^{(j)})\) only the \(j\)-th selector holds, so
\(\mathfrak C\entails\forall x(\delta(x;\hat c^{(j)})\leftrightarrow
\psi_j(x;c))\).

\emph{The rank drops.}  Let \(\theta=\theta_{\delta,k}\) and
\(\operatorname D=\operatorname D_{\delta,k}\), and let \(n=\operatorname
D(\theta)\), a natural number.  As \(\theta\wedge\varphi(x;b)\in p\),
minimality gives \(\operatorname D(\theta\wedge\varphi(x;b))\geq n\).  This
formula implies \(\bigvee_j(\theta\wedge\psi_j(x;c_j))\) in \(\mathfrak C\),
so by items (2) and (4) of \cref{thm:p03-w03-p03-035-drank-calculus} there
is \(j\) with \(\operatorname D(\theta\wedge\psi_j(x;c_j))\geq n\).  Write
\(\theta=\theta(x;a)\) with \(a\) a tuple from \(A\).  For each \(l\),
\(c^l_j\equiv_A c_j\) by \(A\)-indiscernibility, so \(ac^l_j\equiv ac_j\),
and item~(3) gives \(\operatorname D(\theta\wedge\psi_j(x;c^l_j))\geq n\);
by item~(2) and the equivalence above,
\(\operatorname D(\theta\wedge\delta(x;\hat{c}^{l\,(j)}_j))\geq n\), where
\(\hat c^{l\,(j)}_j\) is the coded tuple of \(c^l_j\).  The family
\((\delta(x;\hat c^{l\,(j)}_j):l<\omega)\) is equivalent member by member to
\((\psi_j(x;c^l_j):l<\omega)\), so it is \(k\)-inconsistent.  By the
definition of the rank, \(\operatorname D(\theta)\geq n+1\), contrary to the
choice of \(n\).  Hence no formula of \(p\) forks over \(A\), and \(p\) does
not fork over \(A\).

\emph{Any number of variables.}  Let \(p(x)\) be a complete type over a small
\(B\) in an infinite tuple \(x\) of small length.  For each finite subtuple
\(x_0\) of \(x\), the formulas of \(p\) in the variables \(x_0\) form a
complete type \(p_0\) over \(B\), and the first part gives
\(A_{x_0}\subseteq B\) with \(\abs{A_{x_0}}\leq\abs T\) over which \(p_0\)
does not fork.  Let \(A\) be the union of the \(A_{x_0}\); there are
\(\abs x\) finite subtuples, so \(\abs A\leq\abs T+\abs x\).  A formula of
\(p\) mentions finitely many variables, so it lies in some \(p_0\) (by
\cref{thm:p03-w03-p03-031-saturation}~(4), unused variables do not affect
forking).  If it forked over \(A\), it would imply a finite disjunction of
formulas dividing over \(A\), each of which divides over
\(A_{x_0}\subseteq A\) by \cref{thm:p03-w03-p03-035-dividing-basics}~(2);
so it would fork over \(A_{x_0}\), which is impossible.

\emph{Any model.}  Let \(\mathcal M\entails T\) and \(B\subseteq\dom\mathcal
M\).  The argument above used only (Sat), (Copy) and the completeness of
\(T\), so it runs unchanged in a \(\bar\kappa'\)-saturated elementary
extension \(\mathfrak C'\) of \(\mathcal M\) with
\(\bar\kappa'>\abs{\dom\mathcal M}+\abs x+2^{\abs T}\), which exists under
the Axiom of Choice as recorded in \cref{def:p03-w03-p03-031-monster}; there
both \(B\) and the variable tuple \(x\) are small.  By \cref{thm:p03-w03-p03-031-saturation}~(3), with \(\mathfrak C'\)
in place of \(\mathfrak C\), forking over \(A\) of a formula with parameters
in \(\mathcal M\) is the same whether it is witnessed in an elementary
extension of \(\mathcal M\) or in \(\mathfrak C'\).  So every complete type
over \(B\) consistent with \(\mathcal M\) does not fork over some
\(A\subseteq B\) with \(\abs A\leq\abs T+\abs x\), which is the condition of
\cref{def:p03-simple-theory}.
\end{proof}
\end{annexitem}

\begin{annexitem}{annex:p03-w03-p03-035-ominimal-nip}{thm:p03-w03-p03-035-ominimal-nip}{O-minimal theories are NIP}
\begin{proof}
Say \(c_0<\cdots<c_m\) \emph{alternate in} \(\varphi(x;b)\) when
\(\varphi(c_k;b)\leftrightarrow\neg\varphi(c_{k+1};b)\) for \(k<m\).

\emph{A convex count.}  If \(D\) is a union of \(N\) convex sets (points and
intervals are convex) in a linear order, no \(c_0<\cdots<c_{2N+1}\) alternate
in \(D\).  Otherwise \(N+1\) of these points, those of one parity, lie in
\(D\), and between two consecutive ones lies a point outside \(D\).  Two of the
\(N+1\) points lie in one convex component, and so does every point between
them, a contradiction.

\emph{Uniform bound.}  Suppose that for every \(m\) some model contains \(b\)
and \(c_0<\cdots<c_m\) alternating in \(\varphi(x;b)\).  This is expressed by
a sentence, which then holds in \(\mathfrak C\).  The set of formulas
\(\{u_k<u_{k+1}\}\cup\{\varphi(u_k;y)\leftrightarrow\neg\varphi(u_{k+1};y)\}\)
(\(k<\omega\)) is then finitely realized, and by (Sat) realized by some \(b\)
and \(c_0<c_1<\cdots\).  By the convex count, \(\varphi(\mathfrak C;b)\) is not
a finite union of points and intervals, so \(\mathfrak C\) is not o-minimal,
contrary to hypothesis.  Hence some \(m\) has no alternating
\(c_0<\cdots<c_m\) in \(\mathfrak C\), and by completeness in no model.

\emph{VC-dimension and NIP.}  If a set of \(m+1\) distinct elements, listed
\(c_0<\cdots<c_m\), were shattered, a parameter cutting out
\(\{c_k:k\text{ even}\}\) would make them alternate.  So the VC-dimension of
\(\varphi\) is at most \(m\), in particular finite, and \(\varphi\) is NIP by
\cref{thm:p03-w03-p03-035-ip-characterization}.  As this holds for every
\(\varphi(x;y)\) with \(\abs x=1\), \(T\) is NIP by
\cref{thm:p03-w03-p03-035-one-variable}, (3)\(\Rightarrow\)(1).
\end{proof}
\end{annexitem}

\begin{annexitem}{annex:p03-w03-p03-035-structure-nip}{thm:p03-w03-p03-035-structure-nip}{An o-minimal structure has an NIP theory, given the transfer}
\begin{proof}
The axioms of dense linear orders without endpoints are sentences true in
\(\mathcal M\), so every model of \(\operatorname{Th}(\mathcal M)\) expands a
dense linear order without endpoints.  By the cited
\cref{thm:p03-w03-p03-035-kps}, every model of
\(\operatorname{Th}(\mathcal M)\) has all its parameter-definable unary sets
finite unions of points and intervals, so it is o-minimal.
\Cref{thm:p03-w03-p03-035-ominimal-nip} applies to the complete theory
\(\operatorname{Th}(\mathcal M)\).
\end{proof}
\end{annexitem}

\begin{annexitem}{annex:p03-w03-p03-035-dlo}{thm:p03-w03-p03-035-dlo}{Dense linear orders: o-minimal, NIP, unstable, not simple}
\begin{proof}
\emph{O-minimality.}  Let \(\mathcal M\entails T_{\rm dlo}\), \(\varphi(x,\bar
y)\) a formula with \(x\) a single variable, and \(\bar b\) from
\(\mathcal M\).  By \cref{thm:p03-rn-dlo-pattern-truth}, \(\varphi\) is
equivalent modulo \(T_{\rm dlo}\) to a finite disjunction of order patterns
\(\delta(x,\bar y)\).  For each such \(\delta\), \(\delta(x,\bar b)\) is
empty unless the restriction of \(\delta\) to the \(\bar y\)-slots is the
order pattern of \(\bar b\); if it is, \(\delta\) puts \(x\) either in the
block of some \(y_i\), defining the point \(b_i\), or in a block of its own
between two consecutive values of \(\bar b\), below all of them, above all of
them, or (when \(\bar y\) is empty) alone, defining an open interval, an open
ray, or the whole line.  So \(\varphi(\mathcal M,\bar b)\) is a finite union
of points and intervals.  Models of \(T_{\rm dlo}\) are dense linear orders
without endpoints, so \(T_{\rm dlo}\) is o-minimal in the sense of
\cref{def:p03-o-minimal}.  It is complete by
\cref{thm:p03-rn-dlo-pattern-truth}, so it is NIP by
\cref{thm:p03-w03-p03-035-ominimal-nip}.

\emph{Order property and SOP.}  As \(\mathfrak C\) has no greatest element,
choose \(a_0<a_1<\cdots\) in it.  With \(b_j=a_j\), \(a_i<b_j\) iff \(i<j\), so
\(x<y\) has the order property and \(T_{\rm dlo}\) is unstable.  The sets
\(\{x:x<a_i\}\) increase strictly with \(i\), since \(a_i\) lies in the
\((i+1)\)-st and not in the \(i\)-th, so \(x<y\) has SOP and \(T_{\rm dlo}\) is
not simple by \cref{thm:p03-w03-p03-035-sop-not-simple}.
\end{proof}
\end{annexitem}

\begin{annexitem}{annex:p03-w03-p03-035-rcf}{thm:p03-w03-p03-035-rcf}{Real closed fields: o-minimal, NIP, unstable, not simple}
\begin{proof}
\(\mathsf{RCF}\) is complete (\cref{thm:p03-w03-p03-037-complete}) and every
model is o-minimal (\cref{thm:p03-w03-p03-037-o-minimal}); an ordered field is
dense, since \(x<(x+y)/2<y\) for \(x<y\), and has no endpoints, since
\(x-1<x<x+1\).  So \(\mathsf{RCF}\) is NIP by
\cref{thm:p03-w03-p03-035-ominimal-nip}.  In any model the elements
\(0<1<1+1<\cdots\) are strictly increasing, and as in the proof of
\cref{thm:p03-w03-p03-035-dlo} the formula \(x<y\) has SOP, so
\(\mathsf{RCF}\) is not simple by
\cref{thm:p03-w03-p03-035-sop-not-simple}.
\end{proof}
\end{annexitem}

\begin{annexitem}{annex:p03-w03-p03-035-dense-predicate}{thm:p03-w03-p03-035-dense-predicate}{An NIP ordered theory that is not o-minimal}
\begin{proof}
\emph{(1) One existential quantifier.}  Let \(\gamma(x,\bar y)\) be a
conjunction of literals of \(\Lang_P\) (atomic formulas and their
negations).  We show that \(\exists x\,\gamma\) is equivalent modulo
\(T_{\rm dp}\) to a quantifier-free formula in \(\bar y\).  Literals not
containing \(x\) move outside the quantifier.  Modulo the linear-order
axioms, \(\neg(x=y_i)\) is equivalent to \(x<y_i\vee y_i<x\),
\(\neg(x<y_i)\) to \(x=y_i\vee y_i<x\), and \(\neg(y_i<x)\) to
\(x=y_i\vee x<y_i\); the literals \(x=x\), \(\neg(x<x)\) are equivalent to
\(\top\) and \(x<x\), \(\neg(x=x)\) to \(\bot\).  Rewriting the result in
disjunctive normal form and using that \(\exists x\) distributes over
\(\vee\), it suffices to treat a conjunction whose literals in \(x\) are of
the forms \(x=y_i\), \(y_i<x\) for \(i\in L\), \(x<y_j\) for \(j\in U\),
\(P(x)\), \(\neg P(x)\).  If \(x=y_i\) occurs, \(\exists x\,\gamma\) is
equivalent to \(\gamma\) with \(y_i\) substituted for \(x\).  If both
\(P(x)\) and \(\neg P(x)\) occur, it is equivalent to \(\bot\).  Otherwise
we claim that \(\exists x\,\gamma\) is equivalent to
\(\bigwedge_{i\in L,j\in U}y_i<y_j\) (to \(\top\) when \(L\) or \(U\) is
empty).  One direction is transitivity.  Conversely, let elements \(\bar b\)
of a model satisfy \(b_i<b_j\) for \(i\in L\), \(j\in U\).  If
\(U\neq\varnothing\), let \(c\) be the least \(b_j\), \(j\in U\).  If
\(L\neq\varnothing\), let \(a\) be the largest \(b_i\), \(i\in L\); then
\(a<c\) when \(U\neq\varnothing\).  If \(L=\varnothing\), let \(a\) be an
element below \(c\), or any element when \(U\) is empty too, using that
there is no least element.  If \(U=\varnothing\), let \(c\) be an element
above \(a\), using that there is no greatest element.  Then \(a<c\), every
element strictly between \(a\) and \(c\) lies above every \(b_i\), \(i\in
L\), and below every \(b_j\), \(j\in U\), and the density axioms give \(x\) with \(a<x<c\) and \(P(x)\), and \(x'\) with
\(a<x'<c\) and \(\neg P(x')\); whichever of them matches the literal on
\(P\) satisfies \(\gamma\).

\emph{Quantifier elimination and completeness.}  By induction on formulas,
every formula is equivalent modulo \(T_{\rm dp}\) to a quantifier-free one:
atomic formulas are, the class is closed under \(\neg\) and \(\wedge\), and
for \(\exists x\,\psi\) with \(\psi\) equivalent to a quantifier-free
\(\psi'\), write \(\psi'\) in disjunctive normal form, distribute
\(\exists x\), and apply the previous paragraph to each disjunct; \(\forall\)
is \(\neg\exists\neg\).  A sentence is then equivalent to a quantifier-free
sentence of \(\Lang_P\); there are no constant symbols, so this is a
Boolean combination of \(\top\) and \(\bot\), and its truth value is the
same in every model.  As \(T_{\rm dp}\) has the model
\((\mathbb R,<,\mathbb Q)\), it is complete.

\emph{(2) NIP.}  By (1), every partitioned formula \(\varphi(x;\bar y)\)
with \(\abs x=1\) is equivalent to a Boolean combination of atomic formulas
in \(x,\bar y\), each regarded as a partitioned formula in \((x;\bar y)\);
equivalent formulas shatter the same sets.  Each atomic formula has finite
VC-dimension.  If \(x\) does not occur, the truth value at \((a;\bar b)\)
does not depend on \(a\), so no two-element set is shattered.  If only
\(x\) occurs, as in \(P(x)\), \(x=x\) or \(x<x\), the truth value does not
depend on \(\bar b\), so no one-element set is shattered.  For
\(x<y_i\): if \(a_1<a_2\), no \(\bar b\) has \(a_2<b_i\) and
\(\neg(a_1<b_i)\), so \(\{a_1,a_2\}\) is not shattered; for \(y_i<x\) the
subset \(\{a_1\}\) cannot be cut out in the same way; for \(x=y_i\) the
whole of \(\{a_1,a_2\}\) cannot be cut out.  So each atomic formula is NIP
by \cref{thm:p03-w03-p03-035-ip-characterization}, (1)\(\Leftrightarrow\)(2),
every \(\varphi(x;\bar y)\) with \(\abs x=1\) is NIP by
\cref{thm:p03-w03-p03-035-boolean}, and \(T_{\rm dp}\) is NIP by
\cref{thm:p03-w03-p03-035-one-variable}, (3)\(\Rightarrow\)(1).

\emph{(3) Not o-minimal.}  Let \(\mathcal M\entails T_{\rm dp}\) and suppose
\(P(\mathcal M)\) is a union of finitely many points and intervals.  An
interval \(I\subseteq P(\mathcal M)\) with two elements \(a<b\) contains,
being convex, every \(x\) with \(a<x<b\), and one of these satisfies
\(\neg P(x)\) by the second density axiom; so every interval in the union has
at most one element, and \(P(\mathcal M)\) is finite.  But
\(P(\mathcal M)\) is infinite: given \(a_0<b\), the first density axiom
gives \(a_1\in P(\mathcal M)\) with \(a_0<a_1<b\), then \(a_2\in
P(\mathcal M)\) with \(a_1<a_2<b\), and so on, with all \(a_n\), \(n\geq1\),
distinct.  So \(\mathcal M\) is not o-minimal.

\emph{(4) Order and strict order.}  As \(\mathfrak C\) has no greatest
element, choose \(a_0<a_1<\cdots\) in it.  Then \(a_i<a_j\) iff \(i<j\), the
order property for \(x<y\); and \(\{x:x<a_i\}\subsetneq\{x:x<a_{i+1}\}\)
since \(a_i\) lies in the second set and not in the first, the strict order
property.  So \(T_{\rm dp}\) is unstable, and not simple by
\cref{thm:p03-w03-p03-035-sop-not-simple}.
\end{proof}
\end{annexitem}

\begin{annexitem}{annex:p03-w03-p03-035-random-graph-ip}{thm:p03-w03-p03-035-random-graph-ip}{The edge relation shatters every finite set}
\begin{proof}
Let \(A\) be a finite set of vertices of a model of \(T_{\rm RG}\) and
\(A_0\subseteq A\).  The extension axiom for \(U=A_0\) and
\(V=A\setminus A_0\) gives a vertex \(b\) outside \(A\) adjacent to every
member of \(A_0\) and to none of \(A\setminus A_0\), so
\(\{a\in A:E(a,b)\}=A_0\).  Thus every finite set is shattered by \(E(x;y)\),
item~(2) of \cref{thm:p03-w03-p03-035-ip-characterization} holds, and
\(E(x;y)\) has IP.
\end{proof}
\end{annexitem}

\begin{annexitem}{annex:p03-w03-p03-035-random-graph}{thm:p03-w03-p03-035-random-graph}{The random graph is simple, unstable and has IP}
\begin{proof}
Simplicity and the base: by \cref{thm:p03-w03-p03-031-random-graph}, with
\(C\) the set of entries of \(a\) that lie in \(B\), no entry of \(a\) lies in
\(B\setminus C\), so \(\operatorname{tp}(a/B)\) does not fork over \(C\); and
\(C\) is finite.  Instability: choose distinct vertices
\(a_0,b_0,a_1,b_1,\ldots\) in \(\mathfrak C\) in this order, taking \(a_i\)
adjacent to no earlier \(b_j\) and \(b_j\) adjacent exactly to the earlier
\(a_i\) with \(i<j\), each choice being one application of an extension axiom
to the finite set of earlier vertices.  Then \(E(a_i,b_j)\) iff \(i<j\), the
order property.  The independence property is
\cref{thm:p03-w03-p03-035-random-graph-ip}.
\end{proof}
\end{annexitem}
